\documentclass[]{../../../../tex/classes/styledArticle}
\usepackage{../../../../tex/packages/hebrewSupport}
\usepackage{../../../../tex/packages/mathShortcuts}
\usepackage{../../../../tex/packages/theoremsSupport}

\usepackage{listings}
\lstset{
	language=Python,
	numbers=left,
	breaklines,
	basicstyle=\ttfamily,
	commentstyle=\color{gray},
	prebreak=\textrightarrow,
	tabsize=4
}

\usepackage{halloweenmath}
% lstlisting
\newcommand\ff{\mathrm{f}}

\author{שחר פרץ}
\title{\textit{הרצאה 2} $\sim$ אלגוריתמים}
\begin{document}
	\maketitle
	
	
	\subsection*{תרגילים לשימוש ב־BFS}
	\defi{\textit{קוטר} של עץ זה המרחק הגדול ביותר בין שני צמתים בעץ. }
	\textbf{שאלה: }איך נמצא קוטר של עץ? 
	
	\textbf{הצעה: }נריץ BFS מכל קודקוד, כך נגלה מי הכי רחוק ממנו, ואז ניקח את המקסימום. סיבוכיות: $\oc(\sof V \cdot (\sof V + \sof E))$. ידוע $\sof{E} = \sof{V} - 1$ (כי אנחנו בעץ) וסה''כ קיבלנו $\oc(\sof{V}^{2})$. 
	
	\textbf{רעיון לשיפור: }נריץ BFS מקודקוד כלשהו $s$ בעץ. נסמן ב־$w$ את הקודקוד במרחק המקסימלי מ־$s$. נריץ BFS מ־$w$ ונחזיר את המרחק של הקודקוד הרחוק ביותר מ־$w$. סיבוכיות: $\oc(\sof{V} + \sof{E}) = \oc(\sof{V})$. נוכיח נכונות. \begin{proof}
		ננניח בשלילה שקיימים שני קודקודים $v_1, v_2$ כך ש־$\dg(v_1, v_2) > \dg(w, u)$ לכל $u \in V$. נסמן ב־$v$ את האב הקדמון הראשון של $v_1, v_2$ ו־$w$ (אב קדמון זה דבר מוגדר, כי ברגע שהרצנו BFS קיבלנו עץ מכוון, גם אם התחלנו מגרף לא מכוון). בה''כ $v$ הוא אב קדמון (ראשון) של $v_1$ ו־$w$ (זה אפשרי כי אנחנו בעץ). אז: 
		\[ \dg(v_1, v_2) \le \dg(v_1, v) + \dg(v_2, v) \le \dg(v_1, v) + \dg(v, w) = \dg(v_1, w) \]
		וזו סתירה. 
	\end{proof}
	
	\textbf{שאלה: }נתון גרף $G = (V, E)$ לא מכוון, וצומת $s \in V$. נרצה למצוא לכל $v \in V$ מסלול קצר ביותר באורך \textit{זוגי}. 
	
	\rmark{לפעמים יהיה צורך לחזור על אותה הקשת פעמיים. }
	
	''אני עדיין צריך לעשות $n^{2}$ בדיקות האם אתה שחור``
	
	יש לנו שתי גישות לבעיה הזו. 
	\begin{enumerate}
		\item לשנות את הגרף ולהריץ BFS רגיל. יש לנו כבר הוכחת נכונות ל־BFS, אז יהיה לנו יותר קל להוכיח. 
		\item להריץ BFS עם שינוים. 
	\end{enumerate}
	
	נציג שני פתרונות בשתי הגישות: 
	\begin{enumerate}
		\item נבנה גרף חדש שבו כל קודקוד $v$ משוכפל ל־$v_0, v_1$. הגרף יהיה דו־צדדי. לכל קשת $(u, v)$ בגרף המקורי נוסיף שתי קשתות בגרף החדש, $(u_0, v_1), (u_1, v_0)$. נריץ BFS מ־$s_0$ ונחזיר שהמרחק אל $u$ הוא $\dg(s_0, u_0)$ בגרף החדש. הסיבוכיות: $\oc(\sof V + \sof E)$, כי רק הכפלנו פי 2 את כמות הצמתים והקשתות. נוכיח נכונות. \begin{proof}
			נניח שיש מסלול באורך $2k$ בין $s$ ל־$u$ ב־$G$ המקורי, נראה שניתן לתרגם אותו למסלול באותו האורך ב־$G'$ שבנינו בין $s_0$ ל־$u_0$. בעצם, נחליף את הקודקודים במסלול המקורי בקודקודים ב־$G'$, לסירוגין, בין הצדדים. לדוגמה: 
			\begin{align*}
				s \to a \to b \to c \to d \to e \to u && \mapsto && s_0 \to a_1 \to b_0 \to c_1 \to d_0 \to e_1 \to u_0
			\end{align*}
			
			מצד שני, כל מסלול בין $s_0$ ל־$u_0$ ב־$G'$ בהכרח באורך זוגי, ונוכל לבצע את ההתאמה למסלול שקול ב־$G$ בדומה לכיוון ההפוך. 
		\end{proof}
		\item ניתן לשנות את BFS כך שנתחזק לכל קודקוד $d_0$ ו־$d_1$, שמייצגים את המרחק הקצר ביותר הזוגי והאי־זוגי בהתאמה מ־$s$ אליו. בכל איטרציה נזכור האם זאת איטרציה זוגית או אי־זוגית, ונעדכן בהתאם רק את הערכים הרלוונטיים בהתאמה. 
		כלומר: 
		
		\begin{itemize}
			\item נתחיל באינטרציה $i = 1$. נפצל למקרים: 
			\begin{itemize}
				\item אם $i$ זוגי: עבור צומת $v$ עבורו $\dd_0[v] = i - 1$ ושכן $u$ של $v$ עבורו $\dd_1[v] = \infty$, נעדכן $\dd_1[u] = i$. 
				\item אם $i$ אי־זוגי, עבור צומת $v$ עבורו $\dd_1[v] = i - 1$ ושכן $u$ של $v$ עבורו $\dd_0 = [u]$, נעדכן $\dd_0[u] = i$. 
			\end{itemize}
		\end{itemize}
		נוכיח נכונות: 
		\prooftrivial
		
		סיבוכיות: כל קודקוד יקבל מרחק לכל היותר פעמיים, מרחק זוגי ומרחק אי־זוגי. לכן נעבור על כל קשת לכל היותר 4 פעמים. מכאן שזה לינארי בגודל הגרף. 
	\end{enumerate}
	
	\section*{DFS}
	\textit{חיפוש לעומק}, במקום לרוחב. ב־BFS נסינו למצוא מרחקים מ־$s$. ל־DFS לא איכפת מצומת ספציפית. הוא סורק את כל הגרף, ומוצא את כל רכיבי הקשירות. בניגוד ל־BFS שמאוד בבירור הוא נותן את עץ המסלולים הקצרים ביותר מ־$s$, אין משפט נחמד כזה של מה DFS מחזיר. ''הוא נותן משהו טיפה יותר אמורפי``. 
	
	האלגוריתם תלוי בסדר שבו מייצגים את הקודקודים והשכנים (זה אגב גם נכו ןעל BFS, בהתחלטה מה יופיע או לא יופיע ב־$\pi$). האלגוריתם רקורסיבי. נגדיר את הצבעים הבאים: 
	\begin{itemize}
		\item לבן = צומת שטרם נתקלנו בו
		\item אפור = נתקלנו בו, ולא סיימנו לטפל בו
		\item שחור = סיימנו לטפל בו
	\end{itemize}
	
	(אלו אותם הצבעים מ־BFS). עוד נגדיר את הפרמטרים הבאים: 
	\begin{itemize}
		\item $\pi(u)$ = צומת הקודם (ש''גילה``) את $u$. 
		\item $\dd[u]$ = (discovry) זמן גילוי $u$
		\item $\ff[u]$ = (finish) זמן סיום טיפול ב־$u$
	\end{itemize}
	אכפת לנו רק מיחס הסדר על $\dd, \ff$. אין משמעות למספרים האבסולוטיים עצמם. 
	
	\sen
	\begin{lstlisting}
DFS(G):
	for all u in V: color[u] = white
	time = 0
	for u in V: 
		if u is white: 
			pi[u] = NIL
			DFS-visit(u)

DFS-visit(u): 
	color[u] = grey
	time = time + 1
	d[u] = time
	for v out-neighbor of u: 
		if color[v] = white: 
			pi[v] = u
			DFS-visit(v)
	color[u] = black
	time += 1
	f[u] = time
	\end{lstlisting}
	\she
	
	מה התכונות של זה? הסיבוכיות היא $\oc(\sof V + \sof E)$. אפשר להראות את זה בכך שנראה ש־DFS-visit נקראת רק $\sof V$ פעמים בדיוק (פעם אחת לכל קודקוד). נבחין במספר תכונות: 
	
	\theo{DFS מבקר בכל הצמתים, ומצביעי $\pi$ מייצגים יער (מכוון), אין בו מעגלים (גם במובן הלא מכוון). }\begin{proof}[הסבר]
		כל צומת $v$ בוקר או בלולאה הראשית, ואז $\pi[v] = \mathrm{NIL}$ והוא שורש של אחד העצים. או, בזמן טיפול בקודקוד $u$, ואז $\pi[v] = u$. מכאן ש־$\dd[u] < \dd[v]$. מכאן, אין מעגלים ב־$\pi$, שכן לכל קשת $(v, u)$ ב־$\pi$ התכונה מתקיימת, אז $\pi$ DAG. בנוסף, לכל צומת יש דרגת יציאה של $0$ או $1$. לכן $\pi$ יער. 
	\end{proof}
		
		קשתות $\pi$ מכוונת הפוך מקשתות $G$, והקשתות המתאימות ב־$G$ נקראות \textit{קשתות עץ}. לא כל קשת ב־$G$ תופיע ב־$\pi$. 
		
		יש לנו מיון סוגי קשתות. 
		\begin{enumerate}[i. ]
			\item קשתות עץ – קשתות שמופיעות ב־$\pi$ (בכיוון ההפוך)
			\item קשתות אחוריות/חוזרות/back edges: $(u, v)$ כאשר $v$ אב קדמון של $u$ ביער ה־DFS. בדוגמה לעיל, רק $(x, v)$
			\item קשתות קדמיות/forward edges: קשת $(u, v)$ כאשר $v$ צאצא של $u$ (אבל לא בן ישיר). בדוגמה לעיל,  $(u, x)$. 
			\item קשתות חוצות/cross edges: $(u, v)$ כאשר $v$ לא צאצא ולא קשת של $u$. בדוגמה לעיל, $(w, y)$. 
		\end{enumerate}
		
		את סוגי הקשתות אפשר לאפיין בזמן הריצה של האלגוריתם תוך שימוש בבדיקות על האם הצומת אפור (ולא רק לבן). נראה זאת בהמשך. 
		
	יש לנו תכונה נוספת, שנקראת ''עקרון הקינון``, ונובעת מהמחסנית הרקורסיה. \theo{אם $u, v$ צמתים המקיימים $\dd[u] < \dd[v] < \ff[u]$ אז $v$ צאצא של $u$ ביער ה־DFS, כלומר $\dd[u] < \dd[v] < \ff[v] < \ff[u]$. כלומר, כל זוג אינטרוויאלים $(\dd[v], \ff[v])$ ו־$\dd[u], \ff[v]$, הוא או מקונן או זר. }\begin{proof}
		יהיו $u, v$ המקיימים את התנאים לעיל. אז $v$ התגלה אחרי $u$. בנוסף, $v$ התגלה לפני שסיימנו לטפל ב־$u$. כלומר, $u$ במחסנית הרקורסיה בזמן גילוי $v$. מכיוון שזו מחסנית, נהיה חייבים לסיים ללטפל ב־$v$ לפני סיום הטיפול ב־$u$. 
	\end{proof}
	
	כך נוכל לבצע איפיון קשתות תוך כדי ריצה: 
	\begin{enumerate}[i.]
		\item קשת עץ: $u$ קרא ל־$v$ ברקורסיה
		\item קשת אחרוית: אם $v$ אפור (ולא שחור) כש־$u$ מסתכל על הקשת. מבחינת אינטרוויאלים, $(\dd[u], \ff[u]) \subset (\dd[v], \ff[v])$. \begin{proof}
			מתקיים $\dd[v] < \ff[u]$ כי לפני שגילנו את $u$ גילינו כבר את $v$. מעקרון הקינון, אחד משניים מתקיים: (מימין $u$ התגלה קודם, משמאל $v$ התגלה קודם)
			\begin{align*}
				\dd[v] < \dd[u] < \ff[u] < \ff[v] && \lor && \dd[u] < \dd[v] < \ff[v] < \ff[u]
			\end{align*}
			המצב מימין לא ייתכן, כי $v$ צאצא של $u$ ו־$u$ בודק אותו רק לאחר שהרקורסיה של $v$ חזרה, ואז $v$ אמור להיות שחור והנחנו שהוא אפור. לכן $v$ התגלה קודם, $u$ צאצא של $v$ וזו קשת אחורית. 
		\end{proof}
		\item קשתות חוצות וקדמיות: כאשר $v$ שחור כש־$u$ מסתכל על הקשת, כאשר החלוקה בהתאם ל־: 
		\begin{itemize}
			\item קדמית אם $\dd[u] < \ff[v] < \ff[u]$. במקרה זה האינטרוויאלים מקוננים בכיוון $(\dd[u], \ff[u]) \supset (\dd[v], \ff[v])$. 
			\item חוצה אם: $\ff[v] < \dd[u]$. במקרה זה האינטרוואלים כלומר $(\dd[u], \ff[u]) \cap (\dd[v], \ff[v]) = \varnothing$. ההוכחה בדומה למקודם. 
		\end{itemize}
	\end{enumerate}
	בגרף לא מכוון יש רק קשתות אחריות קשתות עץ (אין כלל קשתות חוצות, ומשום שאנו נסווג קשתות כשנתקל בהן לראשונה, אין קשתות קדמיות). 
	\subsection*{עקרון המסלול הלבן}
	נניח שבגרף $G$ יש מסלול $v_1, v_2 \dots v_k$ כך שב־DFS, $v_1$ הראשון שהתגלה (מבין צמתי המסלול). אזי, כל צמתי המסלול הם צאצאים של $v_1$ ביער ה־DFS. \begin{proof}
		באינדוקציה על $k$. בעבור $k = 1$ טרוויאלי. נניח נכונות על $k - 1$. אזי $v_{k - 1}$ צאצא של $v_1$ ביער. בפרט, $v_{k - 1}$ מסתיים לפני $v_1$. אם (בשלילה) $v_k$ לא צאצא של $v_1$, אז הוא לא אב קדמון, כי אז הוא היה מתגלה קודם. $v_{k - 1}$ מסתיים לפני $v_1$ ולכן כשהוא מסתיים הוא בדק את הקשת $(v_{k - 1}, v_k)$ ומכאן $v_k$ לא לבן בעת זו. 
		
		$v_k$ התגלה אחרי $v_1$ ולפני ש־$v_{k -1}$ הסתיים, ולכן לפני ש־$v_1$ הסתיים. כלומר, $\dd[v_1] < \dd[v_k] < \ff[v_1]$, ולפי עקרון הקינון $v_k$ צאצא של $v_1$. 
	\end{proof}
	האינטואציה של רכיבי קשירות ברורה, אך שימו לב שהיא בעייתית מאוד בגרף מכוון (אפשר לנסות לנסח משהו עם רכיבי קשירות חזקה). 
	
	\subsection*{שימושים של DFS}
	\begin{enumerate}
		\item מציאת רכיבי קשירות בגרף לא מכוון. נריץ DFS, וכל עץ יהיה רכיב קשירות של הגרף. השורש של כל עץ הוא הצומת שנתקלים בו בפונקציה החיצונית. \begin{proof}
			כל עץ ב־DFS קשיר, והוא מורכב מקשתות הגרף, ולכן הוא מוכל ברכיב קשירות במלואו. ולהפך, כל רכיב קשירות נכנס בשלמותו לעץ יחיד ב־DFS לפי עקרון המסלול הלבן. עבור רכיב קשירות $C$, נסמן ב־$u$ את הצומת הראשון שמתגלה. אזי, לכל $v \in C$ יש מסלול לבן מ־$u$ ל־$v$, ולכן $v$ יהיה עם $u$ באותו העץ. 
		\end{proof}
		\item בנוסף, נוכל לשמור לכל קודקוד מי השורש שלו ביער ה־DFS. זה מאפשר לבדוק בקלות האם שני קודקודים נמצאים באותו רכיב הקשירות – פשוט נבדוק האם יש להם את אותו השורש ביער ה־DFS. 
		
		בגרפים מכוונים, יחס השקילות הזה לא מוגדר היטב, אלא משתנה כתלות בסדר מעבר על הקודקודים. 
		\item נתון גרף, מכוון או לא, נרצה לקבוע אם ב־$G$ יש מעגל. 
		
		\theo{יש בגרף מעגל $\iff$ בפלט ה־DFS יש קשת אחורית}
		\begin{proof}
			\begin{itemize}
				\item[$\implies$] אם ב־$G$ יש קשת אחורית $(u, v)$, אז $v$ אב קדמון של $u$ וקשתות העץ מ־$v$ ל־$u$ בתוספת הקשת $(u, v)$ יוצרת מעגל ב־$G$. 
				\item[$\impliedby$]נניח שיש מעגל ב־$G$. נסמן ב־$u$ את הצומת הראשון מהמעגל שהתגלה. נסמן ב־$v$ את הצומת שלפני $u$ במעגל. נניח שמ־$u$ לא הלכנו ישירות ל־$v$. לפי עקרון המסלול הלבן, $v$ יהיה צאצא של $u$ ולכן $(v, u)$ קשת אחורית. 
			\end{itemize}\envendproof
		\end{proof}
		\item מיון טופולוגי של גרף מכוון חסר מעגלים (DAG):
		\defi{\textit{מיון טופולוגי} ב־DAG הוא סידור של הצמתים $v_1, v_2 \dots v_n$ כך שכל קשת $(v_i, v_j)$ מקיימת ש־$i  <j$. }
		
		לא המיון הטופולוגי לא בהכרח יחיד. בדוגמה שאני אעתיק אחכ, אפשר להחליף את הסדר בין 1 ל־2 במיון הטופולוגי. 
		
		נרצה למצוא מיון כזה בעזרת DFS. נוכל פשוט למיין את הקודקודים לפי זמן סיום יורד, וסיימנו. \begin{proof}
			נניח ש־$(u, v)$ קשת  בגרף. צ''ל $f(u) > f(v)$. כלומר ש־$U$ מסתיים אחרי $v$. 
			\begin{enumerate}[A.]
				\item אם $(u, v)$ קשת עץ, זה נובע מיידית מיידית מהרקורסיה. 
				\item אם $(u, v)$ קשת אחורית, אין מעגלים בגרף, ומכאן שאלוהים קיים. 
				\item אם $(u, v)$ קשת קדמית, הכל בסדר כי $\dd[u] < \ff[v] < \ff[u]$. 
				\item אם $(u, v)$ קשת חוצה, אז $\ff[v] < \dd[u] < \ff[u]$ וגם התכונה מתקיימת. 
			\end{enumerate}\envendproof
		\end{proof}
	\end{enumerate}
	
	\defi{\textit{כשר} בגרף לא מכוון הוא קשת שהסרתה מנתקת את רכיב הקשירות שלו. }
	\theo{קשת אינה גשר $\iff$ היא נמצאת על מעגל פשוט בגרף. }\begin{proof}
		\begin{itemize}
			\item[$\implies$]אם $(u, v)$ אינה כשר, אחרי שנורידה מהגרף נקבל מסלול $v \to u$, ובצירוף עם $(u, v)$ נקבל מעגל. 
			\item[$\impliedby$]אם $(u, v)$ על מעגל פשוט, אחרי שנסירה עדיין נוכל לקשר בין שני הקודקודים. 
		\end{itemize}\envendproof
	\end{proof}
	\theo{גרף לא מכוון וקשיר ניתן לכוון לגרף קשיר בחוזקה $\iff$ אין בו קשרים}
	נאמר, שאתם עריית ת''א. מכאן שאתם רוצים להפוך כל רחוב לחד סטרי ובה־בעת שיהיה אפשר להגיע מכל מקום בעיר לכל מקום. אז המשפט הזה הוא בשבילכם. 
	\begin{proof}
		\begin{enumerate}
			\item[$\impliedby$]תהא $(u, ,v)$ קשת, ונניח שכיוונה הוא $u \to v$. כיוון שהגרף קשיר בחוזקה, ישנה קשת מ־$v$ ל־$u$. ביחד עם הקשת נקבל מעגל ולכן הקשת לא כשר, כלומר אין בו גשרים. 
			\item[$\implies$]נתבונן בשטח שת''א ספחה ועכשיו רוצה להפוך אותו לחלק אינטגרלי מהעיר (כלומר בלי נתיבים דו־סטריים). נריץ DFS. נכוון קשתות עץ מאב לבן, קשתות אחוריות מצאצא לאב קדמון. אין עוד סוגי קשתות (אנחנו בגרף קשיר לא מכוון). נניח $s$ הוא שורש של עץ DFS (ויש רק עץ אחד כי הגרף קשיר). 
			
			מספיק להוכיח שניתן להגיע מ־$s$ לכל קודקוד $v$ בשני הכיוונים (ואז כשנרצה מסלול מ־$u$ ל־$v$ נעבור דרך $s$). מסלול $s \xrightflutteringbat{} v$: דרך קשתות עץ. מסלול $v \xrightflutteringbat{} s$: אם $v = s$ סיימנו. אחרת, יהא $u$ הורה של $v$ בעץ ה־DFS. $(u, v)$ איננה גשר, כלומר היא חלק ממעגל. לכן קיימת קשת אחרוית שיוצאת מ־$v$ או צאצא שלו ומגיעה אל $u$ או אב קדמון של $u$. נסמן את הקודקוד שהיא מגיעה אליו ב־$w$. בנינו מסלול מ־$v$ ל־$w$. אם $w = s$ סיימנו, אחרת, נמשיך אינדוקטיבית בטיפוס. בכל טיפוס כזה העומק קטן ממש, ולכן נגיע בסוף ל־$s$. 
		\end{enumerate}\envendproof
	\end{proof}
	לא נספיק היום לשאול את השאלה איזה קודקודים כשמוחקים אותם מקטינים מספר רכיבי קשירות
	
	
	
	
	
	\ndoc
\end{document}